%%
%% This is file `sample-sigconf.tex',
%% generated with the docstrip utility.
%%
%% The original source files were:
%%
%% samples.dtx  (with options: `all,proceedings,bibtex,sigconf')
%% 
%% IMPORTANT NOTICE:
%% 
%% For the copyright see the source file.
%% 
%% Any modified versions of this file must be renamed
%% with new filenames distinct from sample-sigconf.tex.
%% 
%% For distribution of the original source see the terms
%% for copying and modification in the file samples.dtx.
%% 
%% This generated file may be distributed as long as the
%% original source files, as listed above, are part of the
%% same distribution. (The sources need not necessarily be
%% in the same archive or directory.)
%%
%%
%% Commands for TeXCount
%TC:macro \cite [option:text,text]
%TC:macro \citep [option:text,text]
%TC:macro \citet [option:text,text]
%TC:envir table 0 1
%TC:envir table* 0 1
%TC:envir tabular [ignore] word
%TC:envir displaymath 0 word
%TC:envir math 0 word
%TC:envir comment 0 0
%%
%%
%% The first command in your LaTeX source must be the \documentclass
%% command.
%%
%% For submission and review of your manuscript please change the
%% command to \documentclass[manuscript, screen, review]{acmart}.
%%
%% When submitting camera ready or to TAPS, please change the command
%% to \documentclass[sigconf]{acmart} or whichever template is required
%% for your publication.
%%
%%
\PassOptionsToPackage{table}{xcolor} % must precede \documentclass (acmart loads xcolor)
\documentclass[sigconf, screen, review]{acmart}

%% ── Additional packages for GRASP paper ──
\usepackage{booktabs}
\usepackage{multirow}
\let\Bbbk\relax            % avoid clash with newtxmath (loaded by acmart)
\usepackage{amssymb}       % for \checkmark
\usepackage{enumitem}      % for itemize options [nosep, leftmargin=*]
\setlength{\textfloatsep}{10pt plus 2pt minus 2pt} % default 20pt, tighten figure-to-text gap

%%
%% \BibTeX command to typeset BibTeX logo in the docs
\AtBeginDocument{%
  \providecommand\BibTeX{{%
    Bib\TeX}}}

%% Rights management information.  This information is sent to you
%% when you complete the rights form.  These commands have SAMPLE
%% values in them; it is your responsibility as an author to replace
%% the commands and values with those provided to you when you
%% complete the rights form.
\setcopyright{none} 
%\setcopyright{acmlicensed}
\copyrightyear{2026}
\acmYear{2026}
\acmDOI{XXXXXXX.XXXXXXX}

%% These commands are for a PROCEEDINGS abstract or paper.
\acmConference[MICRO 2026]{The 58th IEEE/ACM International Symposium on Microarchitecture}{October 31--November 04, 2026}{Athens, Greece}
%%
%%  Uncomment \acmBooktitle if the title of the proceedings is different
%%  from ``Proceedings of ...''!
%%
%%\acmBooktitle{Woodstock '18: ACM Symposium on Neural Gaze Detection,
%%  June 03--05, 2018, Woodstock, NY}

\acmISBN{978-X-XXXX-XXXX-X/XX/XX}


%%
%% Submission ID.
%% Use this when submitting an article to a sponsored event. You'll
%% receive a unique submission ID from the organizers
%% of the event, and this ID should be used as the parameter to this command.
%%\acmSubmissionID{123-A56-BU3}

%%
%% For managing citations, it is recommended to use bibliography
%% files in BibTeX format.
%%
%% You can then either use BibTeX with the ACM-Reference-Format style,
%% or BibLaTeX with the acmnumeric or acmauthoryear sytles, that include
%% support for advanced citation of software artefact from the
%% biblatex-software package, also separately available on CTAN.
%%
%% Look at the sample-*-biblatex.tex files for templates showcasing
%% the biblatex styles.
%%

%%
%% The majority of ACM publications use numbered citations and
%% references.  The command \citestyle{authoryear} switches to the
%% "author year" style.
%%
%% If you are preparing content for an event
%% sponsored by ACM SIGGRAPH, you must use the "author year" style of
%% citations and references.
%% Uncommenting
%% the next command will enable that style.
%%\citestyle{acmauthoryear}
\settopmatter{printfolios=true}
\settopmatter{printacmref=false}
%\pagestyle{plain}

%%
%% end of the preamble, start of the body of the document source.
\begin{document}

%%
%% The "title" command has an optional parameter,
%% allowing the author to define a "short title" to be used in page headers.
\title{GRASP: Grasping GPU Indirect Memory Patterns via ISA-Exposed Register-Chain Prefetching}
\subtitle{\normalsize{MICRO 2026 Submission
    \textbf{\#NaN} -- Confidential Draft -- Do NOT Distribute!!}}
%%
%% The "author" command and its associated commands are used to define
%% the authors and their affiliations.
%% Of note is the shared affiliation of the first two authors, and the
%% "authornote" and "authornotemark" commands
%% used to denote shared contribution to the research.
%\author{\normalsize{MICRO 2026 Submission
 %   \textbf{\#NaN} -- Confidential Draft -- Do NOT Distribute!!}}

%%
%% By default, the full list of authors will be used in the page
%% headers. Often, this list is too long, and will overlap
%% other information printed in the page headers. This command allows
%% the author to define a more concise list
%% of authors' names for this purpose.

%%
%% The abstract is a short summary of the work to be presented in the
%% article.

%%%%%% -- PAPER CONTENT STARTS-- %%%%%%%%

\begin{abstract}

Indirect memory access (IMA), where one load's address depends on another load's returned data (e.g., $A[B[i]]$), is a major performance bottleneck for graph analytics and sparse linear algebra on GPUs, causing the majority of L1 cache misses. While CPU architects have developed effective IMA prefetchers, their single-threaded assumptions prevent adoption on massively multithreaded GPUs. Meanwhile, existing GPU prefetchers rely on stride-based patterns and cannot handle data-dependent accesses.

We propose GRASP (GPU Register-chain Aware Sector Prefetcher), a hardware prefetcher that exploits a structural property of GPU's SASS ISA: the \texttt{IMAD.WIDE} instruction directly exposes IMA parameters---scale, base address, and index value---as instruction operands, enabling zero-inference detection. GRASP identifies \texttt{LDG}$\to$\texttt{IMAD.WIDE}$\to$\texttt{LDG} dependency chains at the issue stage and employs a two-step index-then-data prefetch pipeline, with per-SM shared tables for training amortization and (base, scale) key merging to handle compiler-induced PC proliferation.

Evaluated on graph workloads with an A100 configuration in Accel-Sim, GRASP achieves a geomean IPC improvement of 35.5\% (up to 49.8\%), at a hardware cost of approximately 2.5\,KB per SM.

\end{abstract}
%%
%% The code below is generated by the tool at http://dl.acm.org/ccs.cfm.
%% Please copy and paste the code instead of the example below.
%%
%\begin{CCSXML}
%<ccs2012>
% <concept>
%  <concept_id>00000000.0000000.0000000</concept_id>
%  <concept_desc>Do Not Use This Code, Generate the Correct Terms for Your Paper</concept_desc>
%  <concept_significance>500</concept_significance>
% </concept>
% <concept>
%  %<concept_id>00000000.00000000.00000000</concept_id>
%  <concept_desc>Do Not Use This Code, Generate the Correct Terms for Your Paper</concept_desc>
%  <concept_significance>300</concept_significance>
% </concept>
% <concept>
%  %<concept_id>00000000.00000000.00000000</concept_id>
%  <concept_desc>Do Not Use This Code, Generate the Correct Terms for Your Paper</concept_desc>
%  <concept_significance>100</concept_significance>
% </concept>
% <concept>
 % <concept_id>00000000.00000000.00000000</concept_id>
%  <concept_desc>Do Not Use This Code, Generate the Correct Terms for Your Paper</concept_desc>
%  <concept_significance>100</concept_significance>
% </concept>
%</ccs2012>
%\end{CCSXML}

%\ccsdesc[500]{Do Not Use This Code~Generate the Correct Terms for Your Paper}
%\ccsdesc[300]{Do Not Use This Code~Generate the Correct Terms for Your Paper}
%\ccsdesc{Do Not Use This Code~Generate the Correct Terms for Your Paper}
%\ccsdesc[100]{Do Not Use This Code~Generate the Correct Terms for Your Paper}

%%
%% Keywords. The author(s) should pick words that accurately describe
%% the work being presented. Separate the keywords with commas.
\keywords{GPU, prefetching, indirect memory access, graph analytics, cache}

\maketitle

\section{Introduction}
\label{sec:intro}

% ¶1: IMA definition + universality
Graph analytics (e.g., social networks, road networks, web graphs),
sparse linear algebra, and data-intensive applications increasingly
rely on GPUs for acceleration~\cite{wang2016gunrock, xu2019gardenia,
yang2022graphblast, chen2020pangolin}.
A defining characteristic of these workloads is the \emph{indirect
memory access} (IMA) pattern \texttt{A[B[i]]}, in which one load's
address depends on the return value of a prior
load~\cite{yu2015imp, fu2024dmp}.
IMA fundamentally breaks spatial and temporal locality: the data
array addresses are determined at runtime and scatter unpredictably
across memory~\cite{talati2021prodigy, xue2024tyche}.
Because this irregularity is inherent to the data structure---not an
artifact of poor layout---it cannot be eliminated by data
reordering~\cite{burtscher2012irregular, shi2018graph}.

% ¶2: IMA breaks GPU latency hiding + Prefetching opportunity + Fig 1
GPUs hide memory latency by interleaving thousands of concurrent
warps, yet IMA undermines this mechanism on two fronts.
First, IMA instructions contribute \textbf{61--93\% of all L1 data cache
misses} across six representative graph algorithms
(Figure~\ref{fig:cache-ima}, left), making them the dominant source
of pipeline stalls.
Second, the serial dependency between index and data loads doubles
the effective memory latency per access, inflating the number of
warps required for full latency hiding well beyond the hardware limit
of 64 warps per SM (\S\ref{sec:motivation}).
%
Crucially, however, a prefetching opportunity exists: while L1 miss
rates are high (38--59\%), L2 hit rates are also high (70--89\%),
indicating that the data \emph{resides in L2 but not in L1}
(Figure~\ref{fig:cache-ima}, left).
An ideal L1 data cache that eliminates all L1 misses achieves a
geometric mean speedup of \textbf{2.04$\times$} across 25~workload
configurations (Figure~\ref{fig:cache-ima}, right), confirming
substantial headroom for an L1-targeted prefetcher.

\begin{figure}[t]
  \centering
  \includegraphics[width=\columnwidth]{figures/fig1_cache_ima_ideal.pdf}
  \caption{(a)~L1 miss rate decomposed into IMA data, IMA index, and
    non-IMA misses (stacked), alongside L2 miss rate (per-algorithm
    median, symmetric graphs).
    (b)~Ideal L1D speedup (geometric mean per algorithm; overall
    2.04$\times$ across 25 configurations).}
  \label{fig:cache-ima}
\end{figure}

% Transition to ¶3
This motivates GPU-specific prefetching for IMA---yet no existing GPU
prefetcher addresses this access pattern.

% ¶3: GPU stride prefetching — 15 years of progress, IMA blind spot
GPU prefetching has been studied for over 15~years, achieving
consistent gains on stride and streaming
patterns~\cite{lee2010manythread, jog2013orchestrated, oh2016apres,
koo2018caps, oh2018wasp, mostofi2023snake}.
These works progressively advanced from inter-thread
prefetching~\cite{lee2010manythread} and scheduling--prefetch
co-design~\cite{jog2013orchestrated}, to CTA-aware leading-warp
detection~\cite{koo2018caps} and variable-length stride
chains~\cite{mostofi2023snake}.
However, IMA requires a fundamentally different detection mechanism:
the data load address is determined by a runtime index value, not by
an arithmetic stride.
Indeed, CAPS explicitly skips data-dependent
accesses~\cite{koo2018caps}, and Snake's chains remain
stride-based~\cite{mostofi2023snake}.

% ¶4: CPU IMA prefetchers — proven effective but cannot migrate
Meanwhile, the CPU community has accumulated nearly a decade of
hardware IMA prefetching, demonstrating that the problem is tractable.
IMP detects \texttt{A[B[i]]} patterns via address
differentials~\cite{yu2015imp}, and the state-of-the-art
DMP~\cite{fu2024dmp} matches four classes of indirect patterns at
0.9\,KB.
Subsequent work has progressively addressed multi-level and
non-linear variants~\cite{cavus2020atp, kaushik2021gretch,
xue2024tyche, talati2021prodigy}.
Yet all these designs assume a single-threaded address stream with
per-core independent state~\cite{ainsworth2019indirect}---assumptions
that collapse under the 64-warp concurrency of a GPU SM,
where na\"{\i}ve per-warp replication would cost
57.6\,KB/SM (\S\ref{sec:motivation}).

% ¶5: Spare Register — sole GPU IMA attempt, decade-long void
The only hardware prefetcher targeting IMA on GPUs is Spare
Register~\cite{lakshminarayana2014spare}, which detects load-pairs
via value-tag propagation and stores prefetched data in idle
registers, achieving 10\% average IPC improvement.
However, it was designed and evaluated on the Fermi
architecture~(2010)---whose cache hierarchy and warp scheduling
differ substantially from modern SM80+
designs.
In the decade since, every subsequent GPU prefetcher has targeted
stride patterns exclusively~\cite{koo2018caps, oh2018wasp,
mostofi2023snake}.

% Transition — zero-position landing
GPU IMA prefetching thus remains an open problem---one that demands a
detection approach native to the GPU's execution model.

% ¶6: IMAD.WIDE — the ISA structural advantage
Rather than migrating CPU techniques, we observe that the GPU's own
ISA provides a more direct path.
Specifically, NVIDIA's SASS instruction \texttt{IMAD.WIDE\;Rd,
Rindex, Rscale, c[base]} encodes \emph{all} the information needed
to prefetch the data load in an IMA chain:
the \emph{base address} of the data array (a kernel-lifetime constant
in constant memory), the \emph{scale factor} (element size, a
compile-time constant), and the \emph{index dependency} (the source
register linking to the preceding index load).
A single observation of this instruction suffices to extract every
prefetch parameter---with \textbf{zero inference}.
This structural guarantee arises because GPU pointers are 64-bit
while graph indices are typically 32-bit integers; the resulting
address computation (32$\times$32$+$64$\to$64) uniquely maps to
\texttt{IMAD.WIDE} in the SASS ISA, making it the compiler's
obligatory choice (\S\ref{sec:motivation}).
Across 12~algorithms compiled for six GPU generations
(Volta through Hopper), 9 of 12 algorithms exhibit $\geq$86\%
\texttt{IMAD.WIDE} conformance, confirming the pattern's
cross-architecture stability.

% ¶7: GRASP overview + selected results
Based on this insight, we propose \textbf{GRASP} (\textbf{G}PU
\textbf{R}egister-chain \textbf{A}ware \textbf{S}ector
\textbf{P}refetcher), which detects
\texttt{LDG$\to$IMAD.WIDE$\to$LDG} register dependency chains at
the issue stage and prefetches through a two-step
index-then-data pipeline.
GRASP's design inherits two structural advantages from
\texttt{IMAD.WIDE}: (1)~\emph{(base,\,scale) merging} collapses
60+ program counters in SpMV into $\sim$2 table entries, eliminating
PC-space explosion; and (2)~\emph{per-SM shared tables} amortize
training across 64~warps, requiring only $\sim$2.3\,KB/SM.
Evaluated on 28~workloads spanning 9~algorithms and 5~datasets,
GRASP achieves a geometric mean IPC improvement of \textbf{18\%}
(up to \textbf{122\%}) over an unmodified A100 configuration.

% ¶8: Contributions
\smallskip\noindent
This paper makes the following contributions:
\begin{itemize}[leftmargin=*, nosep]
\item \textbf{IMA Characterization and IMAD.WIDE Discovery.}
  We present the first systematic quantification of IMA behavior on
  modern GPUs, showing that IMA accounts for 61--93\% of L1 misses
  with an ideal L1D headroom of 2.04$\times$ (geometric mean).
  We identify \texttt{IMAD.WIDE} as a structurally guaranteed ISA
  feature that encodes all IMA prefetch parameters, stable across
  six GPU generations.
\item \textbf{GRASP Design.}
  We propose a zero-inference IMA prefetcher built on
  \texttt{IMAD.WIDE} detection, featuring an index--data two-step
  pipeline, (base,\,scale) merging, and per-SM shared tables at
  $\sim$2.3\,KB/SM---the first dedicated IMA prefetcher for modern
  GPUs.
\item \textbf{Comprehensive Evaluation.}
  We evaluate GRASP on 28~workloads across three SOTA GPU
  prefetcher baselines (Snake, CAPS, Spare Register),
  demonstrating consistent IPC improvements with a geometric mean
  of 18\% (up to 122\%).
\end{itemize}

%%% ====================================================================
%%% §2 Background & Motivation
%%% ====================================================================
\section{Background and Motivation}
\label{sec:motivation}

\subsection{Indirect Memory Access on GPUs}
\textit{[IMA definition, CSR traversal pattern, why stride/stream prefetchers fail.]}

\subsection{Performance Impact of IMA}
\textit{[L1 miss rate 70--77\%, serial dependency chain doubling latency, Little’s Law argument showing warp deficit.]}

\subsection{Why CPU IMA Prefetchers Cannot Migrate}
\textit{[IMP/DMP/Tyche assume single-thread OoO; GPU has 64 warps sharing L1, no per-thread runahead.]}


%%% ====================================================================
%%% §3 Key Insight
%%% ====================================================================
\section{Key Insight: ISA-Exposed IMA Structure}
\label{sec:insight}

\textit{[IMAD.WIDE exposes scale, base, index as instruction operands — zero-inference detection. Structural guarantee: IMAD.WIDE is the only single-instruction 32×32+64→64 in SASS, compiler will not choose 3-instruction alternative. See insight.md Insight~1.]}


%%% ====================================================================
%%% §4 GRASP Design
%%% ====================================================================
\section{GRASP Design}
\label{sec:design}

\subsection{Overview}
\textit{[Architecture figure (Fig.~\ref{fig:grasp-arch}) — CD/CT/TT/IPU/DPU/PRB/ACU/TC in cache hierarchy. Two-step prefetch: index first, then data.]}

% \begin{figure}[t]
%   \centering
%   \includegraphics[width=\columnwidth]{figures/grasp_architecture.pdf}
%   \caption{Overview of GRASP hardware design.}
%   \label{fig:grasp-arch}
% \end{figure}

\subsection{Chain Detection (CD)}
\textit{[Monitor issue stage for LDG→IMAD.WIDE→LDG. Extract (index\_PC, data\_PC, base, scale). Write CT + TT entries.]}

\subsection{Stride Learning (IST)}
\textit{[Per-PC iteration stride tracking. Convergence in 2--3 observations. Cross-kernel persistence.]}

\subsection{Index Prefetch (IPU)}
\textit{[On demand index load: compute prefetch\_addr = addr + stride × distance. Allocate PRB entry.]}

\subsection{Data Prefetch (DPU + ACU)}
\textit{[On index fill: ACU reads returned value, computes data\_addr = base + value × scale. DPU issues data prefetch.]}

\subsection{Throttle Control (TC)}
\textit{[MSHR utilization monitoring. Suppress prefetch when saturated.]}

\subsection{Hardware Budget}

% ============================================================================
% GRASP 论文表格模板集
% ============================================================================
% 风格来源：DMP (HPCA'24), Tyche (TACO'24), Magellan (ISCA'25), Snake (MICRO'23)
%
% 风格规范：
%   - 使用 booktabs 三线表（\toprule / \midrule / \bottomrule），无竖线
%   - 表头加粗居中，数据行文字居中或左对齐，数字右对齐
%   - 表格标题在上方（\caption 在 \begin{tabular} 之前）
%   - "Our work" 行用 \textbf 或 \rowcolor 轻微高亮
%   - 多级表头用 \multicolumn + \cmidrule
%   - 特征矩阵用 \checkmark（需 amssymb 包）
%
% 依赖包（在 main.tex 的 preamble 中添加）：
%   \usepackage{booktabs}
%   \usepackage{multirow}
%   \usepackage{amssymb}       % for \checkmark
%   \usepackage[table]{xcolor} % for \rowcolor (optional)
% ============================================================================


% ────────────────────────────────────────────────────────────────────────────
% Template 1: Simulator Configuration (模拟器参数配置)
% 用途：Methodology 节中描述 baseline GPU 配置
% 参考：DMP TABLE II, Snake TABLE I, Tyche TABLE I
% ────────────────────────────────────────────────────────────────────────────

\begin{table}[t]
  \caption{Baseline GPU configuration.}
  \label{tab:gpu-config}
  \centering
  \small
  \begin{tabular}{@{}ll@{}}
    \toprule
    \textbf{Parameter}          & \textbf{Value} \\
    \midrule
    GPU Model                   & NVIDIA A100 (SM80) \\
    \# SMs                      & 108 \\
    Warp Schedulers / SM        & 4 (GTO) \\
    Max Warps / SM              & 64 \\
    Max CTAs / SM               & 32 \\
    Clock Frequency             & 1.41\,GHz \\
    \midrule
    L1 D-Cache / SM             & 192\,KB, 128B line, 48-way \\
    L1 MSHRs / SM               & 256 entries \\
    L2 Cache (shared)           & 40\,MB, 128B line, 16-way \\
    L2 MSHRs / partition        & 128 entries \\
    \midrule
    DRAM                        & HBM2, 40 channels \\
    DRAM Bandwidth              & 2039\,GB/s (peak) \\
    Interconnect                & Crossbar, 1.41\,GHz \\
    \bottomrule
  \end{tabular}
\end{table}


% ────────────────────────────────────────────────────────────────────────────
% Template 2: Prefetcher Comparison (Baseline 对比配置)
% 用途：列出所有对比的 prefetcher 及其硬件开销
% 参考：DMP TABLE III — "Configuration of Evaluated Prefetchers"
% ────────────────────────────────────────────────────────────────────────────

\begin{table}[t]
  \caption{Configuration of evaluated prefetchers.}
  \label{tab:prefetcher-config}
  \centering
  \small
  \begin{tabular}{@{}lccc@{}}
    \toprule
    \textbf{Prefetcher} & \textbf{Configuration} & \textbf{Category} & \textbf{Storage} \\
    \midrule
    \textbf{GRASP (ours)}
      & \begin{tabular}[c]{@{}c@{}}8-entry CT, 4-entry TT,\\16-entry PRB, $d$=4\end{tabular}
      & HW & 1.2\,KB \\
    \addlinespace
    Stride-INTRA~\cite{}
      & 32-entry stride table / SM
      & HW & 0.5\,KB \\
    \addlinespace
    Spare Reg.~\cite{}
      & 2-step: stride index + pair table
      & SW-HW & 0.8\,KB \\
    \addlinespace
    Snake~\cite{}
      & 16-entry chain table, $L$=4
      & HW & 1.0\,KB \\
    \addlinespace
    No Prefetch
      & ---
      & --- & --- \\
    \bottomrule
  \end{tabular}
\end{table}


% ────────────────────────────────────────────────────────────────────────────
% Template 3: Benchmark / Dataset List (数据集参数)
% 用途：列出评估用的图数据集及其规模
% 参考：DMP TABLE V — "Real-world Datasets for Evaluation"
%       Magellan — benchmark table with Names, Description, #V, #E
% ────────────────────────────────────────────────────────────────────────────

\begin{table}[t]
  \caption{Graph datasets used for evaluation.}
  \label{tab:datasets}
  \centering
  \small
  \begin{tabular}{@{}lrrll@{}}
    \toprule
    \textbf{Dataset} & \textbf{$|V|$} & \textbf{$|E|$ (dir)} & \textbf{$|E|$ (sym)} & \textbf{Domain} \\
    \midrule
    cit-Patents       & 3.77M &  16.5M & $\sim$33M   & Citation  \\
    web-Google        &  916K &   5.1M & $\sim$8.6M  & Web       \\
    flickr            &  821K &   9.8M & $\sim$19.6M & Social    \\
    roadNet-CA        & 1.97M &   2.8M & 2.8M        & Roads     \\
    \midrule
    soc-LiveJournal1  & 4.85M &    69M & $\sim$138M  & Social    \\
    \bottomrule
  \end{tabular}
  \vspace{2pt}
  \parbox{\columnwidth}{\footnotesize
   Top 4 datasets form the primary evaluation; soc-LiveJournal1 is
   used for scalability analysis (BFS, SpMV only).
   All sourced from SNAP~\cite{} / NetworkRepository~\cite{}.
  }
\end{table}


% ────────────────────────────────────────────────────────────────────────────
% Template 4: Kernel Feature Matrix (算法 × IMA 模式 checkmark)
% 用途：展示每个 benchmark 包含哪些 IMA pattern 类型
% 参考：DMP TABLE IV — "Kernel Configurations" with checkmarks
% ────────────────────────────────────────────────────────────────────────────

\begin{table}[t]
  \caption{IMA pattern types present in each benchmark kernel.}
  \label{tab:ima-patterns}
  \centering
  \small
  \begin{tabular}{@{}llrccc@{}}
    \toprule
    \textbf{Benchmark} & \textbf{Category}
      & \textbf{Chains}
      & \multicolumn{3}{c}{\textbf{Datasets}} \\
    \cmidrule(l){4-6}
      & & & \textbf{cit-P} & \textbf{web-G} & \textbf{flickr} \\
    \midrule
    BFS    & Graph traversal  &  7 & \checkmark$^\dagger$ & \checkmark$^\dagger$ & \checkmark$^\dagger$ \\
    SSSP   & Graph traversal  &  7 & \checkmark$^\dagger$ & \checkmark$^\dagger$ & \checkmark$^\dagger$ \\
    BC     & Graph traversal  & 19 & \checkmark$^\dagger$ & \checkmark$^\dagger$ & \checkmark$^\dagger$ \\
    CC     & Graph traversal  &  6 & \checkmark & \checkmark & \checkmark \\
    SpMV   & Sparse LinAlg    & 29 & \checkmark & \checkmark & \checkmark \\
    PR     & Sparse LinAlg    & 57 & \checkmark$^\dagger$ & \checkmark$^\dagger$ & \checkmark$^\dagger$ \\
    VC     & Graph coloring   & 34 & \checkmark & \checkmark & $\times$ \\
    SymGS  & Sparse smoother  &  2 & \checkmark$^\dagger$ & \checkmark$^\dagger$ & $\times$ \\
    \bottomrule
  \end{tabular}
  \vspace{2pt}
  \parbox{\columnwidth}{\footnotesize
   \textbf{Chains}: number of IMA chains detected via \texttt{LDG}$\to$\texttt{IMAD.WIDE}$\to$\texttt{LDG} in SASS.
   $^\dagger$Dual-variant: both directed and symmetric inputs evaluated; best selected.
   $\times$: excluded due to \texttt{MAXCOLOR} overflow.
   All algorithms also run on roadNet-CA (omitted for space).
  }
\end{table}


% ────────────────────────────────────────────────────────────────────────────
% Template 5: GRASP Hardware Budget (硬件预算分解)
% 用途：展示 GRASP 各组件的 entry 数和存储开销
% 参考：DMP TABLE III 的 "Configuration" 列展开为独立表
% ────────────────────────────────────────────────────────────────────────────

\begin{table}[t]
  \caption{GRASP storage budget breakdown (per SM).}
  \label{tab:grasp-budget}
  \centering
  \small
  \begin{tabular}{@{}lcrc@{}}
    \toprule
    \textbf{Component} & \textbf{Entries} & \textbf{Bits/Entry} & \textbf{Total} \\
    \midrule
    CD (Chain Detector)   & 4-deep FIFO & 64  & 32\,B  \\
    CT (Chain Table)      & 8           & 96  & 96\,B  \\
    TT (Target Table)     & 4           & 80  & 40\,B  \\
    IST (Stride Tracker)  & 8           & 48  & 48\,B  \\
    PRB (Prefetch Buffer) & 16          & 128 & 256\,B \\
    ACU                   & 1           & 128 & 16\,B  \\
    TC (Throttle Ctrl)    & 1           & 32  & 4\,B   \\
    \midrule
    \textbf{Total}        &             &     & \textbf{492\,B} \\
    \bottomrule
  \end{tabular}
\end{table}


% ────────────────────────────────────────────────────────────────────────────
% Template 6: Performance Summary (主结果表, full-width)
% 用途：跨所有 workload 汇总 IPC / speedup / prefetch 指标
% 参考：通常作为 full-width table* 跨两栏
% ────────────────────────────────────────────────────────────────────────────

\begin{table*}[t]
  \caption{Performance comparison across all benchmarks (IPC and speedup over no-prefetch baseline).}
  \label{tab:performance}
  \centering
  \small
  \begin{tabular}{@{}l*{5}{r}r@{}}
    \toprule
    & \multicolumn{5}{c}{\textbf{Speedup over No Prefetch}} & \\
    \cmidrule(lr){2-6}
    \textbf{Benchmark}
      & \textbf{Stride} & \textbf{Spare Reg.} & \textbf{Snake} & \textbf{GRASP}
      & \textbf{Ideal L1D} & \textbf{Base IPC} \\
    \midrule
    BFS     & 1.04$\times$ & 1.08$\times$ & 1.06$\times$ & \textbf{1.32$\times$} & 1.75$\times$ & 30.3 \\
    SSSP    & 1.04$\times$ & 1.08$\times$ & 1.05$\times$ & \textbf{1.42$\times$} & 2.46$\times$ & 34.6 \\
    BC      & 1.03$\times$ & 1.06$\times$ & 1.04$\times$ & \textbf{1.25$\times$} & 3.04$\times$ & 44.7 \\
    CC      & 1.00$\times$ & 1.01$\times$ & 1.00$\times$ & \textbf{1.15$\times$} & 2.22$\times$ & --- \\
    SpMV    & 1.00$\times$ & 1.02$\times$ & 1.01$\times$ & \textbf{1.05$\times$} & 2.31$\times$ & 152.0 \\
    PR      & ---          & ---          & ---          & ---                    & ---          & --- \\
    VC      & ---          & ---          & ---          & ---                    & ---          & --- \\
    SymGS   & ---          & ---          & ---          & ---                    & ---          & --- \\
    \midrule
    \textbf{Gmean} & ---$\times$ & ---$\times$ & ---$\times$ & \textbf{---$\times$} & ---$\times$ & --- \\
    \bottomrule
  \end{tabular}
\end{table*}
  % This imports all table templates;
% In the final paper, move individual tables to their appropriate sections.
% For now, they are all included here for preview.


%%% ====================================================================
%%% §5 Implementation
%%% ====================================================================
\section{Implementation}
\label{sec:impl}

\textit{[Accel-Sim / GPGPU-Sim implementation. 8 source files (grasp\_*.h/cc). Trace-driven IMA pair table for runtime value recovery. See 05\_implementation/.]}


%%% ====================================================================
%%% §6 Methodology
%%% ====================================================================
\section{Experimental Methodology}
\label{sec:method}

% Tables 1, 3, 4 belong here in the final paper
\textit{[GPU config (Table~\ref{tab:gpu-config}), datasets (Table~\ref{tab:datasets}), benchmarks (Table~\ref{tab:ima-patterns}), baselines (Table~\ref{tab:prefetcher-config}).]}


%%% ====================================================================
%%% §7 Evaluation
%%% ====================================================================
\section{Evaluation}
\label{sec:eval}

\subsection{Overall Performance}
% Table 6 belongs here
\textit{[IPC speedup across all workloads (Table~\ref{tab:performance}). GRASP vs baselines vs Ideal L1D.]}

\subsection{Prefetch Quality: Accuracy, Coverage, Timeliness}
\textit{[BFS: acc 55.9\%, cov 93.3\%, time 89.6\%. SpMV: acc 28.8\%, cov 15.5\% — ×16 unroll + MSHR saturation.]}

\subsection{Hardware Overhead}
% Table 5 belongs here
\textit{[GRASP budget (Table~\ref{tab:grasp-budget}): ~2.5 KB/SM shared by 64 warps ≈ 39 B/warp.]}

\subsection{Sensitivity Study}
\textit{[PRB entries, CT entries, prefetch degree, lookahead distance. 2×2 panel line plot.]}

\subsection{Ablation Study}
\textit{[Disable CD / IST / DPU / TC individually. SSSP +42.2\% full → breakdown.]}


%%% ====================================================================
%%% §8 Related Work
%%% ====================================================================
\section{Related Work}
\label{sec:related}

\subsection{CPU IMA Prefetchers}
\textit{[IMP, DMP, Tyche, Gretch, Prodigy, Magellan — single-thread assumptions.]}

\subsection{GPU Stride/Stream Prefetchers}
\textit{[MT-Prefetch, APOGEE, Snake, CAPS — stride-only, cannot handle data-dependent addresses.]}

\subsection{GPU Irregular Access Optimizations}
\textit{[Spare Register, DSAP, software prefetching on GPU — limited scope or require compiler support.]}


%%% ====================================================================
%%% §9 Conclusion
%%% ====================================================================
\section{Conclusion}
\label{sec:conclusion}

\textit{[GRASP exploits IMAD.WIDE structural property for zero-inference IMA detection on GPUs. Two-step index-then-data prefetch. Geomean +35.5\% IPC at ~2.5 KB/SM.]}


\begin{acks}
   \textit{[To be added after acceptance.]}
\end{acks}



%%%%%%% -- PAPER CONTENT ENDS -- %%%%%%%%

%%
%% The next two lines define the bibliography style to be used, and
%% the bibliography file.
\bibliographystyle{ACM-Reference-Format}
\bibliography{sample-base}

\end{document}
\endinput
%%
%% End of file `sample-sigconf.tex'.
